\section{Conclusion}

\label{sec:conclusion}

Ringtail provides the first practical lattice-based threshold signature scheme for
blockchain consensus, achieving 0.6-second signing and sub-millisecond verification
with $> 128$-bit classical security. The two-round protocol integrates naturally with
the Lux Quasar dual-certificate mechanism, providing quantum-safe finality anchors
without impacting the BLS fast path. Ringtail's parameters are derived from the
Module-LWE assumption with concrete security estimates, and the scheme is implemented
in production on the Lux network.

\begin{thebibliography}{99}
\bibitem{shor1997polynomial} P. Shor, ``Polynomial-time algorithms for prime factorization and discrete logarithms on a quantum computer,'' \emph{SIAM J. Computing}, 1997.
\bibitem{boldyreva2003threshold} A. Boldyreva, ``Threshold signatures, multisignatures and blind signatures based on the gap-Diffie-Hellman-group signature scheme,'' in \emph{PKC}, 2003.
\bibitem{komlo2020frost} C. Komlo and I. Goldberg, ``FROST: Flexible round-optimized Schnorr threshold signatures,'' in \emph{SAC}, 2020.
\bibitem{canetti2021uc} R. Canetti et al., ``UC non-interactive, proactive, threshold ECDSA with identifiable aborts,'' in \emph{CCS}, 2021.
\bibitem{bendlin2013semi} R. Bendlin, S. Krehbiel, C. Peikert, and A. Rosen, ``Semi-homomorphic encryption and multiparty computation,'' in \emph{EUROCRYPT}, 2013.
\bibitem{boneh2018threshold} D. Boneh et al., ``Threshold cryptosystems from threshold fully homomorphic encryption,'' in \emph{CRYPTO}, 2018.
\bibitem{cozzo2023sharing} D. Cozzo and N. Smart, ``Sharing the LUOV: Threshold post-quantum signatures,'' in \emph{IMA}, 2023.
\bibitem{espitau2024fiat} T. Espitau, M. Tibouchi, and A. Wallet, ``Practical Fiat-Shamir with aborts threshold signatures from lattices,'' 2024.
\bibitem{regev2009lattices} O. Regev, ``On lattices, learning with errors, random linear codes, and cryptography,'' \emph{JACM}, 2009.
\end{thebibliography}

